\documentclass{article}
\usepackage{arxiv}

\usepackage[utf8]{inputenc}
\usepackage[T1]{fontenc}
\usepackage{hyperref}
\usepackage{url}
\usepackage{booktabs}
\usepackage{amsmath, amssymb}
\usepackage{bm}
\usepackage{natbib}
\usepackage{graphicx}

% The emblem lives in the repository-wide assets directory, two levels up.
\graphicspath{{../../assets/}{../../assets/logos/}{../../assets/biota/}{../../assets/diagrams/}}

\title{The Mathematics and Design of WaterStrider}

% arxiv.sty labels every paper "A Preprint" in the running head and under the
% title. This is an internal technical report, not a preprint of a submission.
\renewcommand{\headeright}{Technical Report}
\renewcommand{\undertitle}{Technical Report}

% arxiv.sty's title block sets a 24pt strut above the author names and 0.4in
% below them, which leaves the byline stranded between "Technical Report" and
% the date. This is the same layout with those two gaps halved. Overridden here
% rather than edited in arxiv.sty, which is a vendored file.
\makeatletter
\renewcommand{\@maketitle}{%
  \vbox{%
    \hsize\textwidth
    \linewidth\hsize
    \vskip 0.1in
    \@toptitlebar
    \centering
    {\LARGE\sc \@title\par}
    \@bottomtitlebar
    \textsc{\undertitle}\\
    \vskip 0.1in
    \def\And{%
      \end{tabular}\hfil\linebreak[0]\hfil%
      \begin{tabular}[t]{c}\bf\rule{\z@}{12\p@}\ignorespaces%
    }
    \def\AND{%
      \end{tabular}\hfil\linebreak[4]\hfil%
      \begin{tabular}[t]{c}\bf\rule{\z@}{12\p@}\ignorespaces%
    }
    \begin{tabular}[t]{c}\bf\rule{\z@}{12\p@}\@author\end{tabular}%
  \vskip 0.2in \@minus 0.1in \center{\@date}   \vskip 0.2in
  }
}
\makeatother

\author{
  Daniel Boupha Christensen \\
  Environmental, Health, and Safety, University of California, Berkeley \\
  The Green Initiative Fund \\
}

\begin{document}

% Masthead. Licence at the left, group emblem at the right. This is an internal
% report, so the block carries the licence and the project mark only, and
% deliberately no journal, volume, page range or DOI furniture.
\noindent
\begin{minipage}[c]{0.74\textwidth}
  \footnotesize
  Copyright 2024--2026 Strawberry Creek Monitoring Group contributors.\\
  Licensed under the Apache License, Version 2.0.\\
  \url{https://www.apache.org/licenses/LICENSE-2.0}
\end{minipage}%
\hfill
\begin{minipage}[c]{0.16\textwidth}
  \raggedleft
  \includegraphics[width=\linewidth]{SCMGemblem.png}
\end{minipage}

\maketitle

\begin{abstract}
WaterStrider is a per-sensor statistical detector for Strawberry Creek. It
removes each sensor's normal daily rhythm, runs three complementary tests on
what remains (a rolling Hampel test for spikes, a two-sided cumulative sum
test for shifts that hold, and an exponentially weighted moving average test
for slow drift), and sets every threshold from a tail fitted to clean data at
a false alarm rate chosen in advance. Nothing in the model is trained by
gradient descent. Every number in it is a median, a deviation or a fitted
tail, so each alert can be explained in one sentence. WaterStrider meets the
same detector contract as DuskCrayfish and CobbleShoal, and it also serves as
the simple baseline that the graph models have to beat on the same labelled
events. Here, we set out the mathematics of the residuals, the three tests and
the calibration, the departures from the original plan that the data required,
and what the shipped calibration measured. In brief, the calibration holds its
false alarm budget on held-out months, but at that budget WaterStrider caught
none of the four labelled anomalies.
\end{abstract}

\section{Introduction}
\label{sec:introduction}

The starting point for WaterStrider was Hampel's X84 rule, which flags a value
more than $k$ median absolute deviations (MAD) from the median, with $k = 5.2$
in Hampel's original suggestion~\citep{hampel1985breakdown}. The median and
the MAD are not pulled around by the outliers they are looking for, the rule
needs no training, and its alerts explain themselves. On this creek, however,
the rule alone is a poor detector, for five reasons. (1) A shift that lasts
becomes the new median after half a window, so only its start is flagged,
while a discharge can hold a new level for hours. (2) A slow drift, such as
probe fouling, never stands out from its own window. (3) Storms and daily
cycles look unusual, since rain dilutes conductivity and water temperature
swings every day. (4) A flat signal, such as depth at base flow, has a MAD near
zero, so every small change looks extreme. (5) A fixed $k$ ignores how many
tests run. The plan counted 17 sensors reporting 96 times a day, about
1{,}632 tests a day per rule. Work on automated anomaly detection in river
turbidity, conductivity and level points the same way, and recommends
combining methods to catch more events while keeping false detections
low~\citep{leigh2019framework}.

WaterStrider keeps the Hampel test and adds what each of these limits calls
for: a cumulative sum (CUSUM) test for shifts that hold, an exponentially
weighted moving average (EWMA) test for slow drift, a daily profile and the
rain gate for daily cycles and storms, a floor on every spread for flat
sensors, and a calibrated threshold for each test in place of a fixed $k$. Each
test covers a failure the other two miss. An interquartile range fence or a
plain z-score on the same window would be close to the same test as Hampel's,
and would add false alarms without catching anything new.

In StrawberryWatch, a model is whatever meets the detector contract, and the
contract does not ask for weights. A detector built from medians and fitted
tails is therefore as much a model as DuskCrayfish. WaterStrider also does a
second job. Any graph neural network invites the comparison with a simple
method on the same data, and WaterStrider provides that comparison on the same
labelled events and at the same false alarm budget, which is a fairer
comparison than X84 with a guessed $k$.

\subsection{The calibrated configuration}
\label{sec:configuration}

Table~\ref{tab:config} lists the settings the shipped calibration was fitted
at, read from the model definition and from
\texttt{water\_strider\_calibration.json}, which was made on 14 October 2026.
We refer to these values throughout the mathematics that follows.

\begin{table}[t]
  \centering
  \caption{The calibrated configuration. These are the settings the shipped
    calibration file records, not defaults or recommendations.}
  \label{tab:config}
  \begin{tabular}{ll}
    \toprule
    Setting & Value \\
    \midrule
    Sampling grid & 15 min, 96 steps a day \\
    Level window & 12 hours (48 steps), trailing median \\
    Daily profile & median over the previous 14 days \\
    Share of real readings a window needs & 0.75 \\
    MAD scale & 1.4826 \\
    Spread floor & larger of probe resolution and $0.10 \times$ spread \\
    Hampel window $w$ & 13 steps, including the current one \\
    CUSUM allowance $k$ & 0.5 \\
    EWMA weight $\lambda$ & 0.2 \\
    Gap that resets CUSUM and EWMA & more than 8 missing steps \\
    Declustering gap & 8 steps \\
    Tail start percentiles & 98, then 95, 90, 85, 80 \\
    Minimum clusters for a tail & 20 \\
    Model false alarm rate $q$ & $10^{-4}$ \\
    Rate per test $q / 3$ & $3.33 \times 10^{-5}$ \\
    Calibration span & 4 June 2025 to 29 April 2026 \\
    Split & even months fitted, odd months held out \\
    Sensors calibrated & 30 (10 sites $\times$ 3 variables) \\
    \bottomrule
  \end{tabular}
\end{table}

\section{Data and Notation}
\label{sec:data}

\subsection{Readings on the grid}
\label{sec:grid}

Each sensor is judged on its own readings. We floor every timestamp to the
quarter hour and average the readings that land on the same grid point, so a
sensor becomes one value per 15 min step, with a gap where nothing arrived. We
never fill gaps. DuskCrayfish showed that a filled-in newest reading causes
false alarms, so instead each window is judged only when at least three
quarters of its readings are real. For calibration, we drop readings outside
their plausible range and readings from before the probe's recorded install
date, since those are not that probe's own.

\subsection{Level, daily profile and residuals}
\label{sec:residuals}

The tests run on residuals, the reading with its normal level and daily rhythm
removed. For sensor readings $x_t$, the level $L_t$ is the median of the 48
readings before step $t$, which is 12 hours of history. The daily profile
$P(\cdot)$ assigns a value to each of the 96 quarter hours of the Pacific day.
For each calendar day we take the median of $x_t - L_t$ at each quarter hour,
and the profile in force on a day is the median of those daily values over the
previous 14 days. A quarter hour seen on fewer than three of those days has no
rhythm to remove yet, and its profile value is zero. The residual is then
\begin{equation}
  r_{t} = x_{t} - L_{t} - P\big(\mathrm{slot}(t)\big),
\end{equation}
so a normal afternoon temperature peak scores close to zero. Medians keep a
past event from bending the profile, and the whole computation costs almost
nothing.

The level is a departure from the plan, which subtracted the daily profile
alone. We tried level windows of 6 and 24 hours as well as 12. At 24 hours the
shift test went over budget on the held-out months, and at 6 hours a drift was
absorbed into the level too quickly to be seen.

\subsection{Spread and its floor}
\label{sec:spread}

Each sensor has one robust spread, the scaled MAD of its residuals over the
clean fitted steps,
\begin{equation}
  s = 1.4826 \cdot \operatorname{median}\big|\, r_{t} - \operatorname{median}(r) \,\big| ,
\end{equation}
where the factor makes $s$ approximately one standard deviation for normal
data. Every spread in WaterStrider is a scaled MAD, and the code names say so.
The distinction matters, because 5.2 raw MADs is about 3.5 standard deviations,
while 5.2 scaled MADs is about 5.2.

A steady sensor's MAD can be close to zero, which makes every small change look
extreme. We therefore floor every spread at
\begin{equation}
  f = \max\big(\varepsilon_{v},\ 0.10\, s\big),
\end{equation}
where $\varepsilon_{v}$ is the smallest step the probe reports: 1.0 for
conductivity and depth, 0.1 for temperature and floating conductivity, and
0.01 for dissolved oxygen, each in that probe's reporting units.
Table~\ref{tab:notation} collects the symbols used in the rest of this
document.

\begin{table}[t]
  \centering
  \caption{Notation.}
  \label{tab:notation}
  \begin{tabular}{lll}
    \toprule
    Symbol & Meaning & Dimension \\
    \midrule
    $x_{t}$ & one sensor's reading at step $t$ & scalar \\
    $L_{t}$ & level, 12 hour trailing median & scalar \\
    $P(\cdot)$ & daily profile by quarter hour & 96 \\
    $r_{t}$ & residual & scalar \\
    $s$ & sensor spread, scaled MAD & scalar \\
    $f$ & spread floor & scalar \\
    $z_{t}$ & standardised residual, $r_{t} / \max(s, f)$ & scalar \\
    $H_{t}, C_{t}, E_{t}$ & Hampel, CUSUM and EWMA statistics & scalar \\
    $\rho$ & lag-one autocorrelation of the residuals & scalar \\
    $t_{0}$ & tail start, a percentile of the clean statistic & scalar \\
    $\hat{\sigma}, \hat{\gamma}$ & fitted tail scale and shape & scalar \\
    $n$ & clean fitted steps & scalar \\
    $N_{t}$ & clusters above $t_{0}$ & scalar \\
    $q$ & false alarm rate for the whole model & scalar, $10^{-4}$ \\
    $z_{q}$ & a test's threshold at rate $q/3$ & scalar \\
    \bottomrule
  \end{tabular}
\end{table}

\section{The Three Tests}
\label{sec:tests}

\subsection{Rolling Hampel test}
\label{sec:hampel}

The Hampel statistic is the distance of a residual from the median of its
window, in floored scaled MADs,
\begin{equation}
  H_{t} = \frac{\big|\, r_{t} - \operatorname{median}(\mathcal{W}_{t}) \,\big|}
               {\max\!\big(1.4826 \cdot \operatorname{MAD}(\mathcal{W}_{t}),\ f\big)} ,
\end{equation}
where $\mathcal{W}_{t}$ holds the $w = 13$ residuals ending at step $t$, about
three hours. The window is long enough that one bad reading cannot drag the
median, and short enough that a lasting shift has been absorbed by the time the
CUSUM test has it. The window includes the step itself, so a spike is judged
against the readings around it. This is the moving-window form of the rule
known as the Hampel identifier. The test catches single spikes and sudden
drops, and it misses shifts that last past half the window.

\subsection{Two-sided CUSUM}
\label{sec:cusum}

The CUSUM test accumulates standardised residuals beyond a small allowance,
in both directions~\citep{page1954continuous},
\begin{equation}
  C^{+}_{t} = \max\big(0,\ C^{+}_{t-1} + z_{t} - k\big), \qquad
  C^{-}_{t} = \max\big(0,\ C^{-}_{t-1} - z_{t} - k\big), \qquad
  C_{t} = \max\big(C^{+}_{t}, C^{-}_{t}\big),
\end{equation}
with $k = 0.5$, the usual choice for catching a shift of about one spread. A
shift that holds keeps adding to one of the sums, while a single spike adds
once and decays. The test therefore catches shifts that hold and misses single
spikes by design. Missing steps are skipped, and both sums reset after more
than 8 consecutive missing steps, so a long gap does not carry a sum across it.

We also tried scaling the CUSUM input for slow wander. Residuals here are
strongly autocorrelated, and a running sum of autocorrelated values spreads
wider than one of independent values by the long-run factor
$\sqrt{(1 + \rho) / (1 - \rho)}$. The scaling shrank real shifts below the
allowance, and the shift test then barely fired, so the shipped calibration
has it turned off.

\subsection{EWMA}
\label{sec:ewma}

The EWMA test averages the standardised residuals with geometrically
decreasing weights and reads the average in units of its own spread,
\begin{equation}
  \bar{z}_{t} = \lambda z_{t} + (1 - \lambda)\, \bar{z}_{t-1}, \qquad
  E_{t} = \frac{|\bar{z}_{t}|}{\sqrt{\lambda / (2 - \lambda)}} ,
\end{equation}
with $\lambda = 0.2$, within the range of 0.2 to 0.3 that the NIST handbook
suggests~\citep{nisthandbookewma}. Averaging narrows the spread by
$\sqrt{\lambda / (2 - \lambda)}$, which is why the limits of an EWMA chart are
narrower than those of the raw data. The test catches slow drift and misses
sudden changes. It resets after long gaps in the same way as CUSUM.

\subsection{Combining the tests}
\label{sec:combination}

A step is flagged when any test crosses its own threshold, and the detector
returns one finding per site, sensor and test. CUSUM and EWMA build up over
time, so their statistics already describe changes that last. A Hampel spike,
by contrast, counts only when a second reading agrees, either the step before
or another sensor at the same site. Writing $R_{t} = H_{t} / z_{q}$ for the
spike's share of its own threshold and $R'_{t}$ for the largest share among
the site's other sensors at the same step, the Hampel score that is compared
with the threshold is
\begin{equation}
  \tilde{H}_{t} = z_{q} \cdot \min\!\big(R_{t},\ \max(R_{t-1}, R'_{t})\big).
\end{equation}
A single bad reading is the job of the Settling Pool screen described in the
support modules report, and WaterStrider does not repeat its checks.

\section{Calibration}
\label{sec:calibration}

\subsection{Clean steps and the split}
\label{sec:clean}

We computed every statistic over the full calibration span, with CUSUM and EWMA
carrying their sums as they would in a live check, and fitted only on clean
steps. A step is clean when it lies outside all eight labelled event spans and
the rain gate would not be raising the threshold there. Each sensor needed at
least 14 clean days in the fitted months to be calibrated at all.

We fitted on even months and held out odd months, so that both sides contain
wet and dry weather. The calibration script records that a chronological split
put the rain on the held-out side. Measured over the span, the share of steps
where the rain gate raises the bar is 21.0\% in even months and 10.9\% in odd
months, against 14.1\% and 22.5\% for a chronological split at 1 February 2026,
so both splits leave wet weather on both sides.

\subsection{Declustered tails}
\label{sec:tails}

The residuals are slowly varying. Their lag-one autocorrelation has a median of
0.963 across the 30 calibrated sensors, with a range of 0.701 to 0.993, and
CUSUM and EWMA stay over a line for many steps once they cross it. Steps above a
threshold are therefore far from independent, and a tail fitted to every
exceedance would describe the length of a few runs rather than how often a run
begins. We declustered the statistic, which is the usual remedy when
exceedances are not independent~\citep{coles2001introduction}. Runs over the
starting threshold that lie within 8 steps of each other count as one cluster,
and we fitted the tail to the cluster peaks. In this form $N_{t}$ counts
clusters and $q$ is a rate of alarms per step, not of steps over the line.

For each sensor and test, the tail start $t_{0}$ is the 98th percentile of the
clean fitted statistic, as in the peaks-over-threshold method of
\citet{siffer2017spot}. CUSUM and EWMA can spend the top 2\% of steps in a few
long runs, so where fewer than 20 clusters exceed $t_{0}$ we step down to the
95th, 90th, 85th and 80th percentiles in turn. A Generalised Pareto
distribution with location zero is then fitted by maximum likelihood to the
cluster peaks minus $t_{0}$~\citep{pickands1975gpd}, and the threshold for a
test's share of the budget is
\begin{equation}
  z_{q} = t_{0} + \frac{\hat{\sigma}}{\hat{\gamma}}\left[\left(\frac{(q/3)\, n}{N_{t}}\right)^{-\hat{\gamma}} - 1\right],
  \qquad
  z_{q} = t_{0} - \hat{\sigma} \log\!\left(\frac{(q/3)\, n}{N_{t}}\right) \ \text{ as } \hat{\gamma} \to 0 ,
\end{equation}
using the exponential limit whenever $|\hat{\gamma}| < 10^{-6}$.

A sensor with fewer than 20 clusters even from the 80th percentile shares a
tail with every sensor of its kind. We place the kind's statistics end to end,
with a gap between sensors that no cluster can cross, and fit one tail to the
pooled peaks. The plan fitted a tail per sensor; pooling is a departure that
the thin sensors required.

\subsection{The budget}
\label{sec:budget}

The model's false alarm rate is $q = 10^{-4}$, the same default the support
modules use, split evenly across the three tests so that each test gets
$q/3 = 3.33 \times 10^{-5}$. Splitting by addition is a union bound, so the
real rate of the three tests together is at most $q$~\citep{dunn1961multiple},
and because the tests share residuals it is likely lower.

The calibration file keeps the whole fit for every sensor and test
($t_{0}$, $\hat{\sigma}$, $\hat{\gamma}$, $n$ and $N_{t}$) as well as the
threshold, together with $q$, the span and split used, the settings and a hash
of the fitted data. Keeping the fit, and not only the threshold, is what lets
the sensitivity setting in Section~\ref{sec:sensitivity} move the rate without
recalibrating.

\subsection{The fitted tails}
\label{sec:fitted}

We calibrated 30 sensors, conductivity, depth and temperature at each of 10
sites, and skipped 11. Nine of the skipped sensors had too few readings,
including every sensor at scnf010 and the dissolved oxygen and floating
conductivity probes at botanical\_garden and south\_fork\_0. The remaining two,
dissolved oxygen and floating conductivity at south\_fork\_3, had only 10 and
12 EWMA clusters even from the 80th percentile, and no pooled tail could be
fitted for them. The calibrated sensors had between 16.5 and 88.9 clean days in
the fitted months. Table~\ref{tab:tails} summarises the 90 tails.

\begin{table}[t]
  \centering
  \caption{The 90 fitted tails by test. A pooled tail is shared by every sensor
    of one kind. The median threshold is in the units of each statistic.}
  \label{tab:tails}
  \begin{tabular}{lrrrrrr}
    \toprule
    Test & Own & Pooled & From 98th & $\hat{\gamma} < 0$ & $\hat{\gamma} > 1$ & Median $z_{q}$ \\
    \midrule
    Hampel & 30 & 0  & 28 & 2 & 3  & 38.5 \\
    CUSUM  & 5  & 25 & 0  & 3 & 10 & 3{,}838.0 \\
    EWMA   & 29 & 1  & 1  & 5 & 5  & 113.0 \\
    \bottomrule
  \end{tabular}
\end{table}

The CUSUM tails depend most on pooling. All 10 conductivity sensors and all 10
depth sensors share one CUSUM tail each, as do 5 of the 10 temperature sensors,
and only one EWMA tail is pooled. The pooled conductivity CUSUM threshold is
3{,}838.0 and the pooled depth CUSUM threshold is 221{,}574.1, because the
depth tail's fitted shape is $\hat{\gamma} = 2.77$. A large positive shape
places the threshold far beyond the observed peaks. The same holds for a few
sensors' own tails, such as the kingman\_hall conductivity EWMA
($\hat{\gamma} = 2.00$, $z_{q} = 73{,}312.4$) and the oxford depth Hampel test
($\hat{\gamma} = 2.83$, $z_{q} = 65{,}418.9$).

\subsection{The held-out check}
\label{sec:heldout}

A calibration passes only if the alarm rate on the held-out clean steps comes
within twice the budget. Table~\ref{tab:heldout} gives the result. Every test
passed, and so did the three tests together.

\begin{table}[t]
  \centering
  \caption{Alarms on held-out clean steps (odd months). An alarm is a run over
    the threshold, with runs closer than 8 steps counted as one. The allowance
    is twice the budget.}
  \label{tab:heldout}
  \begin{tabular}{lrrrr}
    \toprule
    Test & Alarms & Steps & Rate & Allowed \\
    \midrule
    Hampel   & 5 & 172{,}713 & $2.89 \times 10^{-5}$ & $6.67 \times 10^{-5}$ \\
    CUSUM    & 5 & 173{,}961 & $2.87 \times 10^{-5}$ & $6.67 \times 10^{-5}$ \\
    EWMA     & 2 & 173{,}961 & $1.15 \times 10^{-5}$ & $6.67 \times 10^{-5}$ \\
    Any test & 9 & 173{,}961 & $5.17 \times 10^{-5}$ & $2.00 \times 10^{-4}$ \\
    \bottomrule
  \end{tabular}
\end{table}

\section{Serving}
\label{sec:serving}

\subsection{Manual sensitivity}
\label{sec:sensitivity}

The calibration sets the default, and a sensitivity setting lets people move
away from it in steps that stay statistically sound. The setting moves the
false alarm rate rather than the thresholds. Each sensor's tail has its own
shape, so scaling every threshold by the same factor would make one sensor much
noisier and barely change another, whereas moving the rate keeps every sensor
equally strict. Each test can be set on its own, and every new threshold comes
from the tail already fitted, so nothing is refitted. Table~\ref{tab:stops}
lists the five stops.

\begin{table}[t]
  \centering
  \caption{The sensitivity stops, as the model's overall rate. Each test gets a
    third. The expected flags assume 1{,}632 tests a day per rule, from the
    plan's count of 17 sensors.}
  \label{tab:stops}
  \begin{tabular}{lrl}
    \toprule
    Stop & Rate $q$ & Expected false flags \\
    \midrule
    Far fewer  & $10^{-5}$ & about one in 2 months \\
    Fewer      & $3 \times 10^{-5}$ & about one in 3 weeks \\
    Calibrated & $10^{-4}$ & about one a week \\
    More       & $3 \times 10^{-4}$ & about one every 2 days \\
    Far more   & $10^{-3}$ & about one and a half a day \\
    \bottomrule
  \end{tabular}
\end{table}

The range stops there for two reasons. Above about 0.02 per test, the threshold
would fall below the 98th percentile that the tail fit starts from, and the
method no longer applies. Well below $10^{-5}$, the threshold lands far beyond
anything in the calibration data, and the number becomes a guess.

A hand-typed threshold $z$ for one sensor and test runs the threshold formula
backwards, to give the rate it implies,
\begin{equation}
  q = \frac{N_{t}}{n}\left(1 + \frac{\hat{\gamma}\,(z - t_{0})}{\hat{\sigma}}\right)^{-1/\hat{\gamma}},
  \qquad
  q = \frac{N_{t}}{n}\, e^{-(z - t_{0}) / \hat{\sigma}} \ \text{ as } \hat{\gamma} \to 0 ,
\end{equation}
and a typed threshold past a negative shape's finite endpoint implies a rate
of zero. At the calibrated stop, the stored threshold is used unchanged.

\subsection{Rain}
\label{sec:rain}

We calibrated on dry-weather steps, and during and after rain the rain gate
raises each threshold by a multiplier, as it does for the other models. The
gate's default rule doubles the bar while it rains and then returns it to the
base value linearly over 36 hours, counting a reading above 0.1 mm as wet.
WaterStrider declares rain as built-in support behaviour, so attaching a
separate rain modulator is refused, since it would raise the thresholds twice.

\subsection{The detector contract}
\label{sec:contract}

WaterStrider reads one JSON file, the calibration, which the existing file
check verifies by hash, and it imports no deep learning library. A check is
medians and running sums over short windows for each sensor, which takes
milliseconds and a few megabytes. Live checks read two days of readings, which
is not enough to rebuild a 14 day profile, so the served detector uses the
latest profile stored in the calibration as a fixed shape and computes the
level from the live readings. Table~\ref{tab:contract} gives its answers to the
rest of the contract.

\begin{table}[t]
  \centering
  \caption{WaterStrider's answers to the detector contract.}
  \label{tab:contract}
  \begin{tabular}{ll}
    \toprule
    Contract piece & WaterStrider \\
    \midrule
    Sites & in-service sites with a calibrated sensor \\
    Rules & \texttt{hampel\_spike}, \texttt{cusum\_shift}, \texttt{ewma\_drift} \\
    Plain wording & a sudden spike, a jump that stayed, a slow drift \\
    Findings & one per site, sensor and rule over its rain-gated line \\
    Steps over the line needed & 1 \\
    Site choices & don't flag, and leave out \\
    Built-in support & rain \\
    Weather & needed while the rain gate is on \\
    Files & \texttt{water\_strider\_calibration.json} \\
    \bottomrule
  \end{tabular}
\end{table}

\section{Evaluation}
\label{sec:evaluation}

We compared WaterStrider with each of its three tests alone, with naive X84 (a
fixed 5.2 raw MADs from a trailing median, with no daily profile and no rain
handling), and with DuskCrayfish and CobbleShoal, each with its shipped
checkpoint and default settings. We count whole events rather than single readings. We ran the
evaluation on 25 October 2026, against the calibration described above.

\subsection{False flags on held-out months}
\label{sec:falseflags}

Table~\ref{tab:flags} gives the false flags per sensor-week on the clean steps
of the held-out months, over 258.9 sensor-weeks. The budget allows about 0.07
a sensor-week. WaterStrider ran at 0.039, about 1.2 a week across the 30
sensors.

\begin{table}[t]
  \centering
  \caption{False flags per sensor-week on held-out clean data, network wide.
    X84 was counted on conductivity sensors only.}
  \label{tab:flags}
  \begin{tabular}{lr}
    \toprule
    Method & Flags per sensor-week \\
    \midrule
    WaterStrider          & 0.039 \\
    Hampel alone          & 0.019 \\
    CUSUM alone           & 0.023 \\
    EWMA alone            & 0.008 \\
    Naive X84, $k = 5.2$   & 5.308 \\
    Naive X84, $k = 347.3$ & 0.409 \\
    \bottomrule
  \end{tabular}
\end{table}

Naive X84 at $k = 5.2$ flagged about 136 times as often as WaterStrider. The
evaluation then searched for the $k$ at which X84 matches WaterStrider's rate,
up to $k = 2{,}000$, and the closest it came was 0.409 a sensor-week at
$k = 347.3$, still about ten times WaterStrider's rate. X84 cannot reach that
rate at any $k$, because on a flat stretch the raw MAD is zero and any change
at all exceeds every multiple of it.

\subsection{Labelled events}
\label{sec:events}

Table~\ref{tab:events} gives which methods flagged conductivity at the labelled
site in each graded event. The four labelled anomalies are the June 2025
discharge on the south fork, both halves of the September 2025 overnight event
on the south fork, and the March 2026 hydrant break into the north fork.

\begin{table}[t]
  \centering
  \caption{Labelled anomalies flagged at the labelled site's conductivity
    sensor. A tick means the method flagged the event; X84 matched is X84 at
    $k = 347.3$.}
  \label{tab:events}
  \begin{tabular}{llcccccc}
    \toprule
    Event & Site & WS & Hampel & CUSUM & EWMA & X84 & X84 matched \\
    \midrule
    June 2025 discharge      & south\_fork\_1 & & & & & \checkmark & \\
    September 2025 overnight & south\_fork\_1 & & & & & \checkmark & \\
    September 2025 overnight & south\_fork\_2 & & & & & \checkmark & \checkmark \\
    March 2026 hydrant       & north\_fork\_0 & & & & & \checkmark & \\
    \midrule
    Caught & & 0 & 0 & 0 & 0 & 4 & 1 \\
    \bottomrule
  \end{tabular}
\end{table}

WaterStrider caught none of the four, and neither did any of its tests alone.
Naive X84 caught all four at $k = 5.2$, at 136 times the false flag rate, and
one at $k = 347.3$. In the same run, DuskCrayfish flagged three of the four
(the June discharge and the September event at south\_fork\_1, and the
hydrant break) and CobbleShoal flagged none.

The events that are not graded give a similar picture. WaterStrider flagged
the November 2025 foam event at north\_fork\_0 through its Hampel test, 40.5
hours after the earliest flag from any method, which was X84's. It did not
flag the June 2025 discharge at south\_fork\_2 or the August 2025 sprinkler
runoff, and no method flagged the botanical garden actuator event. In both
rain-only events, the November 2025 rain and the April 2026 rainfall,
WaterStrider flagged conductivity at one site, codornices, where X84 flagged 7
and 8 sites and DuskCrayfish 3 and 1.

\subsection{Peaks against thresholds}
\label{sec:peaks}

Table~\ref{tab:peaks} shows how close each test came. For every step of the
event window we divided the test's score by its rain-gated line, and the table
gives the largest share each test reached. A value of 1 would be a flag.

\begin{table}[t]
  \centering
  \caption{The largest share of its rain-gated line each test reached at the
    labelled site's conductivity sensor, at the calibrated rate.}
  \label{tab:peaks}
  \begin{tabular}{llrrr}
    \toprule
    Event & Site & Hampel & CUSUM & EWMA \\
    \midrule
    June 2025 discharge      & south\_fork\_1 & 0.062 & 0.172 & 0.081 \\
    September 2025 overnight & south\_fork\_1 & 0.937 & 0.145 & 0.148 \\
    September 2025 overnight & south\_fork\_2 & 0.098 & 0.164 & 0.105 \\
    March 2026 hydrant       & north\_fork\_0 & 0.256 & 0.099 & 0.237 \\
    \bottomrule
  \end{tabular}
\end{table}

Only one event came close. The Hampel test at south\_fork\_1 in the September
event reached 0.937 of its line. Every CUSUM peak stayed below a fifth of the
pooled conductivity threshold of 3{,}838.0, so the shift test, the one meant
for discharges that hold a new level, was furthest from firing.

\subsection{Sensitivity stops}
\label{sec:stopsresult}

We replayed the same events at each sensitivity stop. Table~\ref{tab:stopsresult}
gives the result.

\begin{table}[t]
  \centering
  \caption{Labelled anomalies caught, and sites whose conductivity was flagged
    in the November 2025 and April 2026 rain-only events, at each stop.}
  \label{tab:stopsresult}
  \begin{tabular}{lrrl}
    \toprule
    Stop & Rate $q$ & Anomalies caught & Rain-only sites flagged \\
    \midrule
    Far fewer  & $10^{-5}$          & 0 of 4 & 1 and 1 \\
    Fewer      & $3 \times 10^{-5}$ & 0 of 4 & 1 and 1 \\
    Calibrated & $10^{-4}$          & 0 of 4 & 1 and 1 \\
    More       & $3 \times 10^{-4}$ & 1 of 4 & 1 and 1 \\
    Far more   & $10^{-3}$          & 2 of 4 & 4 and 6 \\
    \bottomrule
  \end{tabular}
\end{table}

At the More stop, the Hampel test caught the September event at
south\_fork\_1. At the Far more stop, the EWMA test also caught the hydrant
break, but the rain-only events then flagged 4 and 6 sites. The flags at
codornices in the rain-only events did not go away at any stop, since their
peaks were several times the line.

\subsection{Pass criteria}
\label{sec:criteria}

The plan set four pass criteria, and the shipped calibration meets one of
them. (1) False flags on held-out clean data came within twice the budget,
which passed. (2) It did not catch the labelled discharge events. (3) It
flagged one site in each of the labelled rain-only events. (4) At the same
false flag rate as naive X84, it caught 0 events against X84's 1.

\section{Limitations and Future Work}
\label{sec:limitations}

The main limitation is that the calibrated thresholds are too far above the
events. WaterStrider holds its false alarm budget on held-out months, but at
that budget it catches none of the four labelled anomalies, and at the loosest
stop it catches two while flagging most sites in a rain-only event. As a
baseline, it therefore does not yet set a meaningful bar for the graph models.
That said, DuskCrayfish flagged three of the same four events in this
evaluation, so the events are detectable from these sensors.

The CUSUM tails are the clearest place to look. Most of them are pooled, the
pooled conductivity threshold sat 5.8 to 10.1 times above the four events'
peaks, and the depth tail's large fitted shape put its threshold at
221{,}574.1. A Generalised Pareto tail with $\hat{\gamma} \ge 1$ has no finite
mean, and a threshold extrapolated from such a fit may say more about the fit
than about the creek. We recommend checking the tail fits against the held-out
peaks, and trying a clamped shape, as CobbleShoal clamps $\gamma$ at zero for
its p-values, before tuning anything else.

The level window may remove part of what the tests are looking for. A 12
hour trailing median takes on a new level once that level has held for 6
hours, half the window, so the residual of a discharge that holds longer
returns toward zero. This may be one reason the shift test stays quiet. We
chose 12 hours over 6 and 24 on the false alarm budget alone, and no event was
part of that choice.

The spike test needs a second reading to agree. A real one-step spike in
conductivity is therefore counted only when the depth or temperature sensor at
the same site spikes at the same step, and otherwise not at all. Settling
Pool, which would take over the job of rejecting single bad readings and so
allow the corroboration rule to be relaxed, is not yet built.

The reason recorded for the blocked split does not hold on the counts in
Section~\ref{sec:clean}. Both splits leave rain on both sides, and the blocked
split puts about twice as many gated steps in the fitted months as in the
held-out ones. The split is still a reasonable one, but it was not chosen for
the reason the calibration script states.

The served detector uses a fixed daily profile from the end of the calibration
span. The plan rebuilt a sensor's profile and thresholds after a replacement,
using the install dates in the inventory. The calibration drops readings from
before a probe's install date, but nothing rebuilds the profile while the
detector runs, so the profile may drift from the creek's rhythm as the
seasons change, until the calibration is rerun.

The expected false flags in Table~\ref{tab:stops} come from the plan's count of
17 sensors, while the shipped calibration judges 30. Across the network, the
calibrated stop therefore allows about two false flags a week rather than one,
and WaterStrider measured about 1.2.

Finally, every number in this document describes the calibration made on 14
October 2026 and the evaluation run on 25 October 2026. The tails and
thresholds are fitted quantities, and they will move when the calibration is
rerun.

\bibliographystyle{unsrtnat}
\bibliography{references}

\end{document}
